\chapter{域扩张与伽罗瓦理论}\label{ch:b3:galois}

每个方程是否都能用根式求解，正如二次方程的求根公式
与卡尔达诺的三次公式所暗示的？能否用直尺和圆规三等分一个角？
这两个问题悬而未决数百年，如今都
被伽罗瓦的一个统一思想——否定地——回答：
为每个多项式配上一个由其根的对称性
构成的有限\emph{群}，再由此群读出答案。本章建立这部“词典”：
\omterm{def:b3:galois:extension}{域扩张}与次数、\omterm{thm:b3:galois:splitting}{分裂域}与
\omterm{def:b3:galois:closure}{代数闭包}、有限域（一个完整的理论——以及
所承诺的 $\mathbb F_q^\times$ 的循环性）、可分性，然后是
\omterm{thm:b3:galois:fundamental}{伽罗瓦对应}本身，附完整证明。我们将收获：
经典作图的不可能性、分圆域的
结构，以及五次方程的根式不可解性——
\cref{ch:b3:groups}中 $A_5$ 的单性
在此大显身手。

\section{扩张、次数与代数性}

\begin{definition}\label{def:b3:galois:extension}
\emph{域扩张}\index{域扩张} $L/K$ 即域
$L$ 包含 $K$ 作为子域；此时 $L$ 是 $K$ 上的向量空间，
而\emph{扩张次数}\index{扩张次数} $[L:K]$ 即
其维数。扩张称为\emph{有限}的，若 $[L:K] <
\infty$。域的\emph{特征}是
$\geq 0$ 的生成元，即 $\Z \to K$，$n \mapsto n\cdot 1$ 之核
的生成元，它为 $0$ 或一个质数 $p$；相应地，$K$ 包含一个最小
子域（\emph{素域}），同构于 $\Q$ 或 $\mathbb
F_p = \Z/p\Z$。
\end{definition}

\begin{theorem}[塔律]\label{thm:b3:galois:tower}
若 $K \subseteq L \subseteq M$，则 $[M:K] = [M:L]\,[L:K]$：若
$(e_i)$ 为 $L$ 在 $K$ 上的基，$(f_j)$ 为 $M$ 在
$L$ 上的基，则 $(e_if_j)$ 为 $M$ 在
$K$ 上的基。\index{塔律}
\end{theorem}

\begin{proof}
生成性：$x \in M$ 可写作 $x = \sum_j \lambda_jf_j$
（$\lambda_j \in L$），而每个 $\lambda_j = \sum_i \mu_{ij}e_i$
（$\mu_{ij} \in K$）：$x = \sum_{i,j}\mu_{ij}e_if_j$。独立性：
$\sum_{i,j}\mu_{ij}e_if_j = 0$ 改写为 $\sum_j (\sum_i
\mu_{ij}e_i)f_j = 0$；内层的和属于 $L$，故为零
（$f_j$ 在 $L$ 上独立）；于是所有 $\mu_{ij} = 0$（$e_i$
在 $K$ 上独立）。
\end{proof}

\begin{definition}\label{def:b3:galois:algebraic}
设 $L/K$，$\alpha \in L$。若存在非零的 $P \in K[X]$ 使得
$P(\alpha) = 0$，则称 $\alpha$ 为\emph{代数元}\index{代数元}
（在 $K$ 上）；$\pi_\alpha$ 是
$\{P : P(\alpha) = 0\}$（$K[X]$ 的\omterm{def:b3:rings:ideal}{理想}）的首一生成元，称为其\emph{极小多项式}\index{极小多项式}，
它是一个不可约多项式
（若 $\pi = QR$ 且 $Q(\alpha) = 0$，则由次数的极小性知 $R$
为常数）。否则称 $\alpha$ 为
\emph{超越的}。记 $K(\alpha)$ 为 $L$ 中
包含 $K$ 与 $\alpha$ 的最小子域，$K[\alpha]$ 为
最小子环。
\end{definition}

\begin{theorem}\label{thm:b3:galois:simple}
若 $\alpha$ 是 $K$ 上的\omterm{def:b3:galois:algebraic}{代数元}，且 $d = \deg\pi_\alpha$，
则
\[
K(\alpha) = K[\alpha] \cong K[X]/(\pi_\alpha),
\qquad [K(\alpha) : K] = d,
\]
基为 $1, \alpha, \dots, \alpha^{d-1}$。反之，若
$[L:K] < \infty$，则每个 $\alpha \in L$ 都是\omterm{def:b3:galois:algebraic}{代数元}，其次数
$\deg \pi_\alpha$ 整除 $[L:K]$。
\end{theorem}

\begin{proof}
求值映射 $K[X] \to L$，$P \mapsto P(\alpha)$，其像为
$K[\alpha]$，核为 $(\pi_\alpha)$；由于 $\pi_\alpha$
不可约，故 $K[X]/(\pi_\alpha)$ 为\emph{域}
（\cref{prop:b3:rings:primemaximal}：在\omterm{def:b3:rings:pidufd}{主理想整环} $K[X]$ 中，
\omterm{def:b3:rings:divisibility}{不可约元}生成\omterm{def:b3:rings:primemaximal}{极大理想}），故 $K[\alpha]$ 是域，
它包含 $K$ 与 $\alpha$，从而等于 $K(\alpha)$。
$1, X, \dots, X^{d-1}$ 的剩余类构成商的一组基（由欧几里得
除法），从而得出基与次数。反之，若 $[L:K]
= n < \infty$：$1, \alpha, \dots, \alpha^n$ 线性相关，
从而得到一个零化多项式；于是 $[K(\alpha):K] =
\deg\pi_\alpha$ 整除 $n$，由\omterm{thm:b3:galois:tower}{塔律}可得。
\end{proof}

\begin{corollary}\label{cor:b3:galois:algclosed}
若 $\alpha, \beta$ 是 $K$ 上的\omterm{def:b3:galois:algebraic}{代数元}，则 $\alpha \pm
\beta$，$\alpha\beta$，$\alpha/\beta$（$\beta \ne 0$）：
$L$ 中在 $K$ 上代数的元素构成 $L$ 的一个子域。
此外，代数性具有传递性：
\omterm{def:b3:galois:algebraic}{代数元}上的\omterm{def:b3:galois:algebraic}{代数元}仍是\omterm{def:b3:galois:algebraic}{代数元}。
\end{corollary}

\begin{proof}
$K(\alpha, \beta) = (K(\alpha))(\beta)$ 在
$K(\alpha)$ 上有限（$\beta$ 在 $K \subseteq K(\alpha)$ 上代数），且
$K(\alpha)/K$ 有限：由\omterm{thm:b3:galois:tower}{塔律}，$[K(\alpha,\beta):K] <
\infty$，因此 $K(\alpha, \beta)$ 的每个元素——包括
所列的四个——都是\omterm{def:b3:galois:algebraic}{代数元}（由\cref{thm:b3:galois:simple}）。
传递性：若 $\beta$ 是 $L$ 上的\omterm{def:b3:galois:algebraic}{代数元}，且 $L/K$
为代数扩张，则 $c_0, \dots, c_{m-1}$ 为
$\pi_{\beta/L}$ 的系数，其生成的域 $F = K(c_0, \dots,
c_{m-1})$ 是 $K$ 的有限扩张（重复运用\omterm{thm:b3:galois:tower}{塔律}），且 $F(\beta)/F$ 为
有限的：$[F(\beta):K] < \infty$，故 $\beta$ 在
$K$ 上代数。
\end{proof}

\begin{example}\label{ex:b3:galois:degrees}
$[\Q(\sqrt2):\Q] = 2$，$[\Q(\sqrt[3]2):\Q] = 3$（$X^3 - 2$ 不可约，
由\omterm{thm:b3:rings:criteria}{艾森斯坦判别法}）；$[\Q(\zeta_p):\Q] = p - 1$，其中
$\zeta_p = \eu^{2\iu\pi/p}$（$\Phi_p$ 不可约，
见\cref{ex:b3:rings:eisenstein}）。\omterm{thm:b3:galois:tower}{塔律}在此已是
一件有力的武器：$\sqrt[3]2 \notin \Q(\sqrt2)$，因为 $3 \nmid 2$。
\end{example}

\section{分裂域；代数闭包}

\begin{theorem}[分裂域]\label{thm:b3:galois:splitting}
设 $P \in K[X]$ 为非常数多项式。存在 $P$ 在 $K$ 上的一个\emph{分裂
域}\index{分裂域}：扩张 $L =
K(\alpha_1, \dots, \alpha_n)$，由 $P$ 的根生成，且在其中
$P$ 分解为一次因式之积。它在
$K$-同构意义下唯一，且 $[L:K] \leq (\deg P)!$。
\end{theorem}

\begin{proof}
\emph{存在性}，对 $\deg P$ 作归纳：取一个不可约
因式 $Q$（$P$ 的因式）；域 $K_1 = K[X]/(Q)$ 含有
根 $\alpha_1 = \bar X$，它是 $Q$ 的根，从而也是 $P$ 的根；写 $P = (X -
\alpha_1)P_1$ 在 $K_1$ 上，对 $P_1$ 在 $K_1$ 上应用归纳假设；
于是次数之积至多为 $n(n-1)\cdots = n!$。

\emph{唯一性}由更强的\emph{同构扩张
引理}推出：设 $\sigma \colon K \to K'$ 为同构，
$P \in K[X]$，$P^\sigma$ 为系数经映射所得的多项式，
$L, L'$ 分别为 $P, P^\sigma$ 的\omterm{thm:b3:galois:splitting}{分裂域}，则 $\sigma$ 可延拓
为同构 $L \to L'$。对 $[L:K]$ 作归纳：若 $P$ 在
$K$ 中分裂，则 $L = K$，且 $L' = K'$（$P^\sigma$ 在 $K'$ 中分裂，
且 $L'$ 由其根生成）。否则，
取 $\alpha \in L \setminus K$，它为不可约因式 $Q$（$P$ 的因式）的根，
且 $\deg Q \geq 2$；$Q^\sigma$ 是 $P^\sigma$ 的不可约因式，
有根 $\beta \in L'$，于是
\[
K(\alpha) \cong K[X]/(Q) \xrightarrow{\ \sigma\ }
K'[X]/(Q^\sigma) \cong K'(\beta)
\]
这是 $\sigma$ 的延拓，且 $\alpha \mapsto \beta$。现在 $L$ 是
$P$ 在 $K(\alpha)$ 上的\omterm{thm:b3:galois:splitting}{分裂域}，而 $L'$ 是 $P^\sigma$ 在
$K'(\beta)$ 上的\omterm{thm:b3:galois:splitting}{分裂域}，且 $[L : K(\alpha)] < [L:K]$：
归纳即得 $L \to L'$ 的延拓。
\end{proof}

\begin{definition}\label{def:b3:galois:closure}
称域 $\Omega$ 为\emph{代数闭}\index{代数闭}域，
即 $\Omega[X]$ 中每个非常数多项式都
在 $\Omega$ 中有根（从而分裂）。$K$ 的\emph{代数
闭包}\index{代数闭包}是代数扩张 $\bar K/K$，
且 $\bar K$ 代数闭。
\end{definition}

\begin{theorem}[施泰尼茨]\label{thm:b3:galois:steinitz}
每个域 $K$ 都有\omterm{def:b3:galois:closure}{代数闭包}，且在
$K$-同构意义下唯一。
\end{theorem}

\begin{proof}
\emph{存在性（阿廷构造）。}设 $R = K[(X_f)_f]$ 为
关于变元 $X_f$ 的多项式环，每个非常数
首一多项式 $f \in K[X]$ 对应一个变元，$I$ 为由全体
$f(X_f)$ 生成的\omterm{def:b3:rings:ideal}{理想}。$I$ 是真\omterm{def:b3:rings:ideal}{理想}：一个关系 $1 = \sum_{i=1}^r g_i\,
f_i(X_{f_i})$ 只涉及有限个多项式；在
$E$（$f_1\cdots f_r$ 的公共\omterm{thm:b3:galois:splitting}{分裂域}）中取 $\alpha_i$ 为 $f_i$ 的根，
代入 $X_{f_i} \mapsto \alpha_i$（其余变元
$\mapsto 0$）：$1 = 0$，矛盾。设 $\mathfrak m \supseteq I$ 为
\omterm{def:b3:rings:primemaximal}{极大理想}（见\cref{thm:b3:rings:krull}；佐恩引理），令 $K_1 = R/\mathfrak
m$：它是 $K$ 的\omterm{def:b3:galois:extension}{域扩张}，其中每个非常数 $f \in
K[X]$ 都有根，根为 $\bar X_f$，且
它在 $K$ 上代数（它由各个 $\bar X_f$ 生成，每个都代数）。
迭代：$K \subseteq K_1 \subseteq K_2 \subseteq \cdots$，其中
$K_{n+1}$ 对 $K_n$ 所做的，与 $K_1$ 对 $K$ 所做的相同，设 $\Omega =
\bigcup_n K_n$，它是域。任一非常数 $g \in \Omega[X]$ 的
有限个系数都位于某个 $K_n$ 中；
$g$ 在 $K_n$ 上的不可约因式有根，该根在 $K_{n+1} \subseteq
\Omega$ 中，故 $\Omega$ \omterm{def:b3:galois:closure}{代数闭}，且在
$K$ 上代数（各 $K_n$ 由传递性都代数，
见\cref{cor:b3:galois:algclosed}）：$\Omega$ 是一个\omterm{def:b3:galois:closure}{代数
闭包}。

\emph{唯一性。}设 $\Omega, \Omega'$ 为两个\omterm{def:b3:galois:closure}{代数
闭包}。考虑有序对 $(E, \tau)$ 构成的集合，其中 $K
\subseteq E \subseteq \Omega$，$\tau\colon E \to \Omega'$ 是
$K$-嵌入，按延拓排序；它非空（$(K,
\mathrm{id})$），且满足归纳条件（链的并），由佐恩引理得
极大元 $(E_0, \tau_0)$。若 $E_0 \neq \Omega$，取 $\alpha \in
\Omega\setminus E_0$：$\pi_{\alpha/E_0}$ 映为一个
$\tau_0(E_0)$ 上的多项式，它有根 $\beta$（在\omterm{def:b3:galois:closure}{代数闭}的
$\Omega'$ 中）。$\tau_0$ 可延拓为 $E_0(\alpha) \to
\Omega'$（$\alpha \mapsto \beta$），与极大性矛盾。故
存在 $K$-嵌入 $\tau \colon \Omega \to \Omega'$；其
像与 $\Omega$ 同构，是\omterm{def:b3:galois:closure}{代数闭}的，且
$\Omega'$ 在其上代数：对 $x \in \Omega'$，
$\pi_{x/\tau(\Omega)}$ 在 $\tau(\Omega)$ 上分裂，故 $x \in
\tau(\Omega)$。于是 $\tau$ 为满射，从而是同构。
\end{proof}

\begin{remark}\label{rem:b3:galois:closureexamples}
当 $K = \Q$ 时，可避开超限方法：
代数数 $\bar\Q = \{z \in \C : z \text{ 在
} \Q\}$ 上代数的数构成一个\omterm{def:b3:galois:closure}{代数闭包}——它是 $\C$ 的子域（由
\cref{cor:b3:galois:algclosed}），它\omterm{def:b3:galois:closure}{代数闭}是因为 $\C$
\omterm{def:b3:galois:closure}{代数闭}（达朗贝尔–高斯定理，在
\cref{ch:b3:holomorphic}中用复分析证明），且 $\bar\Q$ 上多项式的根
在 $\Q$ 上代数（由传递性）。
\end{remark}

\section{有限域}

\begin{theorem}\label{thm:b3:galois:finitefields}
设 $p$ 为质数，$n \geq 1$，$q = p^n$。
\begin{enumerate}
  \item 有限域的元素个数必为素数幂，且
        对每个 $q$，恰存在一个域 $\mathbb F_q$，
        它有 $q$ 个元素，在同构意义下唯一，即
        $X^q - X$ 在 $\mathbb F_p$ 上的\omterm{thm:b3:galois:splitting}{分裂域}。\index{有限域}
  \item \emph{弗罗贝尼乌斯自同态}\index{弗罗贝尼乌斯自同态}：$F
        \colon x \mapsto x^p$ 是 $\mathbb
        F_q$ 的自同构，而 $\mathbb F_q$ 的自同构群
        是阶为 $n$ 的循环群，由 $F$ 生成。
  \item $\mathbb F_{p^m}$ 可嵌入 $\mathbb F_{p^n}$，当且仅当 $m \mid
        n$。
\end{enumerate}
\end{theorem}

\begin{proof}
(1) 有限域 $E$ 的特征为 $p > 0$，且是
$\mathbb F_p$ 上的有限维向量空间：$\abs E = p^n$。
其乘法群的阶为 $q - 1$，故每个 $x \in E$
满足 $x^q = x$：$E$ 由 $q$ 个 $X^q - X$ 的根组成，
因此它是其在 $\mathbb F_p$ 上的\omterm{thm:b3:galois:splitting}{分裂域}——
由此至同构意义下确定 $E$
（\cref{thm:b3:galois:splitting}）。反之，设 $L$ 为
$X^q - X$ 的\omterm{thm:b3:galois:splitting}{分裂域}，其根的集合 $E$ 是一个
\emph{子域}：$(x + y)^q = x^q + y^q$，反复应用
“新生之梦” $(a+b)^p = a^p + b^p$（$p \mid \binom pk$），以及
$(xy)^q = x^qy^q$，$(x^{-1})^q = (x^q)^{-1}$。它恰有 $q$
个元素，因为 $X^q - X$ 可分：其导数为
$qX^{q-1} - 1 = -1$（因 $p \mid q$），与之互素，故无
重根。于是 $L = E$ 恰有 $q$ 个元素。

(2) $F$ 是域同态（由新生之梦），单射
（域之间的同态），因而在有限域 $\mathbb F_q$ 上为双射。$F^n =
\mathrm{id}$（由 $x^q = x$），且更小的幂都不是恒等映射：
若 $F^m = \mathrm{id}$，则全部 $q$ 个元素都是
$X^{p^m} - X$ 的根，从而 $p^m \geq q$。故 $\langle F\rangle$
是 $n$ 阶循环群；且再无其他自同构，由
下文证明的界 $\abs{\operatorname{Aut}} \leq [\,\mathbb F_q :
\mathbb F_p\,] = n$ 可知
（\cref{prop:b3:galois:embeddings}，取 $L = \mathbb F_q$，$K =
\mathbb F_p$；自同构固定素域）。

(3) 若 $\mathbb F_{p^m} \subseteq \mathbb F_{p^n}$，由塔式律
得 $p^n = (p^m)^d$，故 $m \mid n$。反之若 $m \mid n$，则
$p^m - 1 \mid p^n - 1$（等比求和），故 $X^{p^m} - X$ 整除
$X^{p^n} - X$（对指数作同样的论证：$X^{a} - 1 \mid X^{b}
- 1$（当 $a \mid b$ 时），前者的根在
$\mathbb F_{p^n}$ 中构成所求子域，其基数为
$p^m$（可分性同 (1)）。
\end{proof}

\begin{theorem}[循环性]\label{thm:b3:galois:cyclic}
域的乘法群的每个有限子群都是
循环群。特别地，$\mathbb F_q^\times \cong
\Z/(q-1)\Z$。\index{有限域的乘法群}
\end{theorem}

\begin{proof}
设 $G \leq K^\times$ 为有限群。由结构定理
（\cref{cor:b3:modules:abelian}），$G \cong \Z/d_1 \times \dots
\times \Z/d_s$，其中 $d_1 \mid \dots \mid d_s$。每个 $x \in G$
则满足 $x^{d_s} = 1$；而 $X^{d_s} - 1$ 至多有
$d_s$ 个根（在域 $K$ 中），故 $\abs G = d_1\cdots d_s \leq d_s$，
于是 $s = 1$，$G$ 是循环群。
\end{proof}

\begin{example}\label{ex:b3:galois:f8}
$\mathbb F_8 = \mathbb F_2[X]/(X^3 + X + 1)$：这个三次式在
$\mathbb F_2$ 中没有根，故不可约。记 $\omega =
\bar X$：$\mathbb F_8^\times$ 是 $7$ 阶循环群，故
\emph{每个} $\neq 0, 1$ 的元素都是生成元。
$\mathbb F_{p^{12}}$ 的子域构成 $12$ 的因数格：
$\mathbb F_p, \mathbb F_{p^2}, \mathbb F_{p^3}, \mathbb F_{p^4},
\mathbb F_{p^6}, \mathbb F_{p^{12}}$——这是第一个完整的
\omterm{thm:b3:galois:fundamental}{伽罗瓦对应}的实例。
\end{example}

\section{可分性与嵌入}

\begin{definition}\label{def:b3:galois:separable}
称 $P \in K[X]$ 为\emph{可分多项式}\index{可分多项式
}，若它在\omterm{thm:b3:galois:splitting}{分裂域}中没有重根——
等价地 $\gcd(P, P') = 1$（重根是公共根；
反之在\omterm{thm:b3:galois:splitting}{分裂域}上，公共根是
重根；且最大公因式在\omterm{def:b3:galois:extension}{域扩张}下不变，
由 \cref{cor:b3:modules:descent} 的论证）。\omterm{def:b3:galois:algebraic}{代数元}
若其\omterm{def:b3:galois:algebraic}{极小多项式}可分，则称其可分；\omterm{def:b3:galois:extension}{域扩张} $L/K$ 若
其全部元素都可分。
\end{definition}

\begin{proposition}\label{prop:b3:galois:perfect}
\emph{不可约}多项式 $P \in K[X]$ 可分，除非 $P' = 0$，
这迫使 $\operatorname{char} K = p > 0$ 且 $P \in K[X^p]$。
因此，特征为 $0$ 的域的每个代数扩张，
以及有限域的每个代数扩张，都是可分的（这样的域称为
\emph{完全域}\index{完全域}）。
\end{proposition}

\begin{proof}
$\gcd(P, P')$ 整除 $P$；若它不为 $1$，由不可约性
知 $\gcd = P$（相差一个常数），从而 $P \mid P'$，且 $\deg
P' < \deg P$，故 $P' = 0$。设 $P = \sum a_kX^k$，$ka_k = 0$
对所有 $k$ 成立，因而特征为 $0$ 时 $P$ 为常数（已排除）；
在特征 $p$ 中，$a_k = 0$ 除非 $p \mid k$：$P =
Q(X^p)$。在有限域上，每个元素都是 $p$ 次幂
（\omterm{thm:b3:galois:finitefields}{弗罗贝尼乌斯自同态}是满射），故 $Q(X^p) = \sum b_k^p X^{pk} = (\sum
b_kX^k)^p$ 不是不可约的，这说明 $P' = 0$ 在有限域上
对不可约的 $P$ 同样不可能出现。
\end{proof}

\begin{proposition}[嵌入计数]\label{prop:b3:galois:embeddings}
设 $L = K(\alpha_1, \dots, \alpha_r)$ 为 $K$ 上的有限扩张，且
$\sigma \colon K \to \Omega$ 为到一个\omterm{def:b3:galois:closure}{代数闭}域中的嵌入，
则 $\sigma$ 到 $L$ 的延拓个数
至多为 $[L:K]$，若 $L/K$ 可分则取等号。
特别地，$\abs{\operatorname{Aut}_K(L)} \leq [L:K]$。
\end{proposition}

\begin{proof}
对 $[L:K]$ 作归纳，借助单步扩张进行。对 $L = K(\alpha)$：一个
延拓 $\tau$ 由 $\tau(\alpha)$ 决定，它必须是
$\Omega$ 中 $\pi_\alpha^\sigma$ 的根；反之，每个这样的
根都给出一个延拓（$K(\alpha) \cong K[X]/(\pi_\alpha)$）。
延拓的个数即 $\pi_\alpha^\sigma$ 在 $\Omega$ 中的\emph{互异}
根的个数，至多为 $\deg\pi_\alpha =
[K(\alpha):K]$，当且仅当 $\pi_\alpha$ 可分时取等号
（$\pi^\sigma$ 与 $\pi$ 的可分性相同：与导数的最大公因式
被 $\sigma$ 保持）。一般地，把 $L =
K(\alpha_1)(\alpha_2, \dots)$ 分解：$\sigma$ 到
$K(\alpha_1)$ 的延拓个数 $\leq [K(\alpha_1):K]$，而每个延拓
可以 $\leq [L : K(\alpha_1)]$ 种方式继续延拓（由归纳法），再相乘（由塔式
律）。在可分情形下，两个计数都取等号：
在更大的域 $K(\alpha_1)$ 上的\omterm{def:b3:galois:algebraic}{极小多项式}整除
$K$ 上的相应多项式，因而仍然可分。
\end{proof}

\begin{theorem}[本原元定理]\label{thm:b3:galois:primitive}
每个有限\emph{可分}扩张都是单扩张：$L =
K(\gamma)$，对某个 $\gamma$ 成立。\index{本原元定理}
\end{theorem}

\begin{proof}
若 $K$ 为有限域，则 $L$ 亦为有限域，
$\gamma$ 为循环群 $L^\times$ 的生成元（\cref{thm:b3:galois:cyclic}）即满足要求。
设 $K$ 为无限域。由归纳法，只需处理 $L =
K(\alpha, \beta)$ 的情形。令 $n = [L:K]$。由
\cref{prop:b3:galois:embeddings}，存在 $n$ 个互异的
$K$-嵌入 $\sigma_1, \dots, \sigma_n \colon L \to \Omega$
（$\Omega$ 为一\omterm{def:b3:galois:closure}{代数闭包}）。多项式
\[
D(T)=\prod_{i<j}\bigl[\bigl(\sigma_i(\alpha)-\sigma_j(\alpha)
\bigr) + T\bigl(\sigma_i(\beta) - \sigma_j(\beta)\bigr)\bigr]
\]
不恒为零：某个因子恒为零，仅当
$\sigma_i, \sigma_j$ 在 $\alpha$ 与 $\beta$ 上皆相同，从而在
$L$ 上相同——与 $i \neq j$ 矛盾。由于 $K$ 为无限域，取 $c \in
K$ 使 $D(c) \ne 0$：于是 $n$ 个元素 $\sigma_i(\alpha +
c\beta)$ 两两不同，故 $\gamma = \alpha + c\beta$
有至少 $n$ 个不同的共轭（均在 $\Omega$ 中），即
$\deg \pi_\gamma \geq n$：$[K(\gamma):K] \geq n = [L:K]$ 迫使
$L = K(\gamma)$。
\end{proof}

\section{伽罗瓦对应}

\begin{definition}\label{def:b3:galois:galois}
有限扩张 $L/K$ 称为\emph{伽罗瓦扩张}\index{伽罗瓦扩张}，
若它是 $K$ 上某个\emph{可分}多项式的\omterm{thm:b3:galois:splitting}{分裂域}。
其\emph{伽罗瓦群}\index{伽罗瓦群}为
$\operatorname{Gal}(L/K) = \operatorname{Aut}_K(L)$，即
$L$ 的域自同构中逐点保持 $K$ 中元素不动者所成的群。
\end{definition}

\begin{proposition}\label{prop:b3:galois:galoisorder}
若 $L/K$ 为\omterm{def:b3:galois:galois}{伽罗瓦扩张}，则 $\abs{\operatorname{Gal}(L/K)} =
[L:K]$；此外，$L/F$ 对每个中间域 $K
\subseteq F \subseteq L$ 都是\omterm{def:b3:galois:galois}{伽罗瓦扩张}，且每个 $F$-嵌入 $L \to \Omega
\supseteq L$ 的像都是 $L$（\emph{正规性}）。
\end{proposition}

\begin{proof}
设 $L$ 为\omterm{def:b3:galois:separable}{可分多项式} $P$ 在 $K$ 上的\omterm{thm:b3:galois:splitting}{分裂域}，并固定一个\omterm{def:b3:galois:closure}{代数闭包}
$\Omega \supseteq L$。$L/K$ 可分：它由
$P$ 的根生成；每个元素的可分性可由
下面的等号情形推出，但我们直接论证——
将\cref{prop:b3:galois:embeddings}应用于生成元（
即可分的 $P$ 的根，其\omterm{def:b3:galois:algebraic}{极小多项式}整除 $P$），
得到恰好 $[L:K]$ 个 $K \hookrightarrow \Omega$ 的延拓
（归纳步中，$P$ 的根在某个中间域上的\omterm{def:b3:galois:algebraic}{极小多项式}
仍整除 $P$，从而可分）。每个这样的嵌入
$\tau \colon L \to \Omega$ 置换 $P$ 的根，
（$\tau$ 固定系数），而 $L$ 由这些根生成，故 $\tau(L) = L$。于是嵌入
$=$ 自同构；
即 $\abs{\operatorname{Gal}(L/K)} = [L:K]$。对中间域 $F$，
$L$ 也是 $P$ 在 $F$ 上的\omterm{thm:b3:galois:splitting}{分裂域}，
且 $P$ 仍可分：$L/F$ 为\omterm{def:b3:galois:galois}{伽罗瓦扩张}；
同样的论证给出 $F$ 上的正规性。
\end{proof}

\begin{lemma}[阿廷]\label{lem:b3:galois:artin}
设 $G$ 为域 $L$ 的有限自同构群，并令
$K = L^G = \{x : \sigma(x) = x\ \forall\sigma \in G\}$ 为其
不动域。则 $[L : L^G] \leq \abs G$。
\end{lemma}

\begin{proof}
设 $n = \abs G$，$G = \{\sigma_1, \dots, \sigma_n\}$，
并假设 $x_1, \dots, x_{n+1} \in L$ 在
$K$ 上线性无关。关于 $n$ 个方程、$n+1$ 个未知量的齐次线性方程组，
未知量为 $(c_j)$，系数取自 $L$，
\[
\sum_{j=1}^{n+1} c_j\,\sigma_i(x_j) = 0
\qquad (i = 1, \dots, n),
\]
有非零解。取一个非零分量个数\emph{最少}的解，
不妨（重新编号）设其非零分量为 $c_1, \dots, c_r \neq 0$，$r \geq 2$
（单个 $c_j\sigma_i(x_j) = 0$ 不可能），并规范化为
$c_r = 1$。并非所有 $c_j$ 都在 $K$ 中：对 $\sigma_i =
\mathrm{id}$ 的方程将与线性无关性矛盾；不妨设 $c_1 \notin K$，
故 $\tau(c_1) \ne c_1$，即存在 $\tau \in G$。对所有方程作用 $\tau$：
由于 $\tau\sigma_i$ 遍历 $G$，向量
$(\tau(c_j))_j$ 亦为一个解；相减得 $(c_j -
\tau(c_j))_j$ 为一个解，且非零分量更少
（第 $r$ 个分量 $1 - 1 = 0$ 为零，第一个分量则非零），且不为零：矛盾。
因此任意 $n+1$ 个元素线性相关：$[L:K]
\leq n$。
\end{proof}

\begin{theorem}[伽罗瓦理论基本定理]
\label{thm:b3:galois:fundamental}
设 $L/K$ 为\omterm{def:b3:galois:galois}{伽罗瓦扩张}，其群 $G =
\operatorname{Gal}(L/K)$。\index{伽罗瓦对应}
\begin{enumerate}
  \item $L^G = K$。
  \item 映射 $H \mapsto L^H$ 与 $F \mapsto
        \operatorname{Gal}(L/F)$ 互为逆映射，
        是 $G$ 的子群与
        中间域 $K \subseteq F \subseteq L$ 之间包含关系反转的双射；
        此外 $[L : L^H] = \abs H$，且 $[L^H : K] = [G : H]$。
  \item $H \trianglelefteq G$ 当且仅当 $L^H/K$ 为\omterm{def:b3:galois:galois}{伽罗瓦扩张}，此时
        限制映射诱导 $\operatorname{Gal}(L^H/K) \cong
        G/H$。
\end{enumerate}
\end{theorem}

\begin{proof}
(1) 显然 $K \subseteq L^G$。反之，设 $\alpha \in L
\setminus K$，我们构造 $\sigma \in G$，使 $\sigma(\alpha) \ne
\alpha$。$\pi_\alpha$ 为 $K$ 上的\omterm{def:b3:galois:algebraic}{极小多项式}，其次数
$\geq 2$，且可分（$L/K$ 可分，
\cref{prop:b3:galois:galoisorder}），故它在一\omterm{def:b3:galois:closure}{代数闭包}中
另有一根 $\beta \neq \alpha$（\omterm{def:b3:galois:closure}{代数闭包}取作 $\Omega \supseteq L$）。
将 $K$-嵌入 $K(\alpha) \to \Omega$，$\alpha \mapsto
\beta$，延拓为嵌入 $\tau\colon L \to \Omega$
（\cref{prop:b3:galois:embeddings}）；由正规性
（\cref{prop:b3:galois:galoisorder}）得 $\tau(L) = L$，故 $\tau \in
G$，$\beta = \tau(\alpha) \in L$，且 $\tau(\alpha) \neq
\alpha$。

(2) 对子群 $H$：$L/L^H$ 是\omterm{def:b3:galois:galois}{伽罗瓦扩张}
（\cref{prop:b3:galois:galoisorder}），且
$\operatorname{Gal}(L/L^H) \supseteq H$ 显然成立，故 $[L:L^H] =
\abs{\operatorname{Gal}(L/L^H)} \geq \abs H$；阿廷引理
给出 $[L:L^H] \leq \abs H$，故等号成立，并且
$\operatorname{Gal}(L/L^H) = H$。对中间域 $F$，
$L/F$ 为\omterm{def:b3:galois:galois}{伽罗瓦扩张}，故 $L^{\operatorname{Gal}(L/F)} = F$，由 (1)
对 $L/F$ 应用得之。两映射互逆；它们显然反转包含关系。
次数方面：刚证得 $[L:L^H] = \abs H$，并且
$[L^H:K] = [L:K]/[L:L^H] = \abs G/\abs H$。

(3) 对 $\sigma \in G$ 与 $H \leq G$：$\sigma(L^H) =
L^{\sigma H\sigma^{-1}}$（直接验证）。由双射，
$\sigma(L^H) = L^H$ 对所有 $\sigma$ 成立，当且仅当 $H \trianglelefteq G$。
现设 $H \trianglelefteq G$，取 $F = L^H$：每个 $\sigma \in
G$ 都限制为 $F$ 的自同构，由此得到同态 $\rho
\colon G \to \operatorname{Aut}_K(F)$，其核为 $\{\sigma :
\sigma\restriction_F = \mathrm{id}\} = \operatorname{Gal}(L/F) =
H$，故 $G/H$ 嵌入 $\operatorname{Aut}_K(F)$，于是
$\abs{\operatorname{Aut}_K(F)} \geq [G:H] = [F:K]$；反向不等式
总成立（\cref{prop:b3:galois:embeddings}）：
$\abs{\operatorname{Aut}_K(F)} = [F:K]$，且 $\rho$ 为满射。
还需证明 $F/K$ 为\omterm{def:b3:galois:galois}{伽罗瓦扩张}：$F$ 关于 $K$ 可分
（位于可分的 $L/K$ 之内），且 $F = K(\gamma)$
（\cref{thm:b3:galois:primitive}）；多项式
$\prod_{\sigma \in G/H}\bigl(X - \sigma(\gamma)\bigr)$（对互异的像求积，
这些像都属于 $F$：$\sigma(F) =
L^{\sigma H\sigma^{-1}} = L^H = F$，由 $H$ 的正规性）
其系数
被 $G$ 固定，由 (1) 知其属于 $K$：它是 $K[X]$ 中的\omterm{def:b3:galois:separable}{可分多项式}，
在 $F$ 上分裂，其根生成 $F$：$F/K$
为\omterm{def:b3:galois:galois}{伽罗瓦扩张}。反之，若 $F = L^H$ 且 $F/K$ 为\omterm{def:b3:galois:galois}{伽罗瓦扩张}，
$F$ 的正规性（\cref{prop:b3:galois:galoisorder}，应用于
嵌入 $F \to \Omega$，由 $G$ 中元素限制而得）给出
$\sigma(F) = F$ 对所有 $\sigma \in G$ 成立，即 $H
\trianglelefteq G$。
\end{proof}

\begin{omfigure}
\begin{tikzpicture}[xscale=1.5, yscale=1.15]
  \tikzset{fd/.style={font=\small}}
  % subgroup lattice of S3 (left, upside down to mirror fields)
  \node[fd] (G)  at (-3.2, 0)   {$S_3$};
  \node[fd] (A)  at (-4.4, 1.2) {$\langle(1\,2\,3)\rangle$};
  \node[fd] (T1) at (-3.6, 1.2) {$\langle(1\,2)\rangle$};
  \node[fd] (T2) at (-2.8, 1.2) {$\langle(1\,3)\rangle$};
  \node[fd] (T3) at (-2.0, 1.2) {$\langle(2\,3)\rangle$};
  \node[fd] (E)  at (-3.2, 2.4) {$\{e\}$};
  \draw (G) -- (A); \draw (G) -- (T1); \draw (G) -- (T2);
  \draw (G) -- (T3);
  \draw (E) -- (A); \draw (E) -- (T1); \draw (E) -- (T2);
  \draw (E) -- (T3);
  % field lattice (right)
  \node[fd] (Q)  at (1.6, 0)   {$\Q$};
  \node[fd] (W)  at (0.4, 1.2) {$\Q(\iu\sqrt3)$};
  \node[fd] (R1) at (1.4, 1.2) {$\Q(\sqrt[3]2)$};
  \node[fd] (R2) at (2.4, 1.2) {$\Q(j\sqrt[3]2)$};
  \node[fd] (R3) at (3.5, 1.2) {$\Q(j^2\sqrt[3]2)$};
  \node[fd] (L)  at (1.6, 2.4) {$L = \Q(\sqrt[3]2, j)$};
  \draw (Q) -- (W); \draw (Q) -- (R1); \draw (Q) -- (R2);
  \draw (Q) -- (R3);
  \draw (L) -- (W); \draw (L) -- (R1); \draw (L) -- (R2);
  \draw (L) -- (R3);
  \draw[->, omDef, thick, shorten >=6pt, shorten <=6pt]
    (-1.4, 1.2) -- (-0.6, 1.2);
\end{tikzpicture}

{\small $L$ 为 $X^3 - 2$ 在 $\Q$ 上的\omterm{thm:b3:galois:splitting}{分裂域}，其\omterm{thm:b3:galois:fundamental}{伽罗瓦对应}
（$j = \eu^{2\iu\pi/3}$）：
$\operatorname{Gal}(L/\Q) \cong S_3$ 的子群（左侧，阶次反序）
与中间域（右侧）相对应。唯一的正规真
子群 $\langle(1\,2\,3)\rangle$ 对应唯一的子扩张
$\Q(\iu\sqrt3)/\Q$，它是\omterm{def:b3:galois:galois}{伽罗瓦扩张}；三个共轭子群
$\langle(i\,j)\rangle$ 分别对应三个共轭三次域
$\Q(j^k\sqrt[3]2)$，它们均
不在 $\Q$ 上正规。}
\end{omfigure}

\section{分圆扩张}

\begin{definition}\label{def:b3:galois:cyclotomic}
设 $n \geq 1$，$\zeta_n = \eu^{2\iu\pi/n}$。
称 $n$ 次\emph{分圆多项式}\index{分圆多项式}为
$\Phi_n = \prod_{\gcd(k,n)=1,\ 1 \le k \le n} \bigl(X -
\zeta_n^k\bigr)$，次数为 $\varphi(n)$；按精确阶对
$X^n - 1$ 的根分组，得 $X^n - 1 = \prod_{d \mid n}\Phi_d$，
由此归纳可知 $\Phi_n \in \Z[X]$（用首一整系数多项式的
带余除法）。
\end{definition}

\begin{theorem}\label{thm:b3:galois:cyclotomicirred}
$\Phi_n$ 在 $\Q$ 上不可约，故 $[\Q(\zeta_n) : \Q] =
\varphi(n)$，并且
\[
\operatorname{Gal}\bigl(\Q(\zeta_n)/\Q\bigr) \;\cong\;
(\Z/n\Z)^\times,
\qquad \sigma_a(\zeta_n) = \zeta_n^a .
\]
因此扩张 $\Q(\zeta_n)/\Q$ 为\omterm{def:b3:galois:galois}{伽罗瓦扩张}，
且其\omterm{def:b3:galois:galois}{伽罗瓦群}为\emph{阿贝尔}群。
\end{theorem}

\begin{proof}
设 $f = \pi_{\zeta_n}$，则 $\Phi_n = fg$，其中 $f, g \in \Z[X]$
首一（高斯引理，见\cref{lem:b3:rings:gauss}：\omterm{def:b3:rings:content}{容度}
相乘，所有多项式均首一）。\emph{断言：若 $\zeta$ 为
$f$ 的根，$p \nmid n$ 为素数，则 $\zeta^p$ 为
$f$ 的根。} 否则 $\zeta^p$ 为 $g$ 的根（它是本原的
$n$ 次单位根），故 $\zeta$ 为 $g(X^p)$ 的根，且 $f
\mid g(X^p)$ 在 $\Z[X]$ 中成立（由\omterm{def:b3:galois:algebraic}{极小多项式}，再次应用高斯引理）。
模 $p$ 约化：$\bar g(X^p) = \bar g(X)^p$（\omterm{thm:b3:galois:finitefields}{弗罗贝尼乌斯
自同态}在 $\mathbb F_p[X]$ 上逐系数有 $a^p = a$，并利用
新生之梦），故 $\bar f \mid \bar g^{\,p}$：$\bar f$ 与 $\bar g$
有一个公共的不可约因子，且 $\bar\Phi_n = \bar f\bar g$ 有
重因子。于是 $X^n - \bar 1$ 也有重因子；但其
导数 $\bar nX^{n-1}$ 与它互素（$p \nmid n$，且
$0$ 不是根），矛盾。

每个本原根 $\zeta_n^k$（$\gcd(k, n) = 1$）都由
$\zeta_n$ 经若干次不整除 $n$ 的素数幂（因子 $k$）得到：
该论断依次传递，故每个本原根都是
$f$ 的根：$f = \Phi_n$，不可约。于是
$[\Q(\zeta_n):\Q] = \varphi(n)$，且 $\Q(\zeta_n)$ 是
\omterm{def:b3:galois:separable}{可分多项式} $X^n - 1$ 的\omterm{thm:b3:galois:splitting}{分裂域}（其根皆为
$\zeta_n$ 的幂）：\omterm{def:b3:galois:galois}{伽罗瓦扩张}。自同构 $\sigma$ 将 $\zeta_n$
映为另一本原根 $\zeta_n^{a(\sigma)}$，且 $\sigma
\mapsto a(\sigma)$ 是到
$(\Z/n\Z)^\times$ 的单态射，两群阶皆为 $\varphi(n)$：
同构。
\end{proof}

\section{尺规作图}

\begin{definition}\label{def:b3:galois:constructible}
将平面等同于 $\C$，从 $\{0, 1\}$ 出发。称一点为
\emph{可构造的}\index{可构造的}，若它可由有限次下述操作的交点得到：
过两个已作出的点作直线，
以已作出的点为圆心、以两个已作出的点之间的距离为半径作圆，
并取这些直线与圆的交点。
\end{definition}

\begin{theorem}[万策尔]\label{thm:b3:galois:wantzel}
$z \in \C$ 可构造，当且仅当存在域塔 $\Q = F_0
\subseteq F_1 \subseteq \dots \subseteq F_r$，使得 $[F_{i+1} :
F_i] = 2$ 且 $z \in F_r$。特别地，可构造数
是次数为 $2$ 的幂的
$\Q$ 上的代数数。\index{万策尔定理}
\end{theorem}

\begin{proof}
($\Rightarrow$) 两直线交点的坐标
由经过坐标属于子域 $F \subseteq \R$ 的点的两直线确定，
满足 $F$ 上的线性方程组，故仍在 $F$ 中。直线与圆、
圆与圆的交点，在消去线性部分后
（两圆方程相减得一直线），
化为 $F$ 上的二次方程：新坐标在 $F$ 中，
或在 $F(\sqrt d)$ 中，$d \in F$，$d > 0$。由归纳法，
每个已作出点的坐标都位于
$\Q$ 的二次扩张塔中；$z = x + \iu y$ 亦在二次塔中
（添加 $\iu$ 即多一步二次扩张）。
次数的推论：$[\Q(z):\Q]$ 整除 $[F_r : \Q] = 2^r$
（由\omterm{thm:b3:galois:tower}{塔律}）。

($\Leftarrow$) 可构造数构成一个域：和与
差可用平行四边形作出（平行线可作图：
两次作垂线即可——经典的
过一点作垂线的方法只需一个圆和两段弧）；
积与商则由泰勒斯截线构形给出，
（给定长度 $a, b$，作出 $ab$ 与
$a/b$，借助两条射线上的相似三角形）。域对
平方根封闭：对 $a > 0$，以 $1 + a$ 为直径的圆与
接点处的垂线相交于高 $\sqrt a$ 处
（直角三角形中高的比例中项关系）；对复数
$w = \rho\eu^{\iu\theta}$，作出 $\sqrt\rho$，
并二等分 $\theta$（二等分角是尺规作图）。
因此二次塔中元素的实部与虚部皆可构造，
由塔上的归纳法得出：每一步添加
一个二次方程的根，它可由域运算
与一个已作出数的平方根表示，
（即二次求根公式；特征为 $0$）。
\end{proof}

\begin{corollary}\label{cor:b3:galois:impossible}
下列三个经典问题都无法仅用直尺和
圆规完成：\index{尺规作图}
\begin{enumerate}
  \item \emph{倍立方问题}：$\sqrt[3]2$ 的次数为
        $3$，不是 $2$ 的幂。
  \item \emph{三等分角问题}：三等分 $60^\circ$
        需要 $\cos 20^\circ$，它是不可约多项式 $8X^3
        - 6X - 1$ 的根，次数为 $3$。
  \item \emph{化圆为方问题}：$\sqrt\pi$ 是超越的
        （$\pi$ 是——林德曼定理，在本阶段不加证明地承认：其
        证明属于超越数论课程）。
\end{enumerate}
此外，正 $n$ 边形可构造，当且仅当
$\varphi(n)$ 为 $2$ 的幂（高斯–万策尔；“若”部分用到
周末问题的方法，“仅当”部分为
\cref{thm:b3:galois:wantzel} 应用于 $\zeta_n$，其次数为
$\varphi(n)$）。对 $n = 7$，$\varphi(7) = 6$，正
七边形不可作图；对 $n = 17$，$\varphi(17) = 16 = 2^4$，
可构造——周末问题即作出了它。
\end{corollary}

\begin{proof}
(1) $X^3 - 2$ 不可约（\omterm{thm:b3:rings:criteria}{艾森斯坦判别法}）。(2) 由 $\cos
3\theta = 4\cos^3\theta - 3\cos\theta$，取 $3\theta =
60^\circ$：$8c^3 - 6c = 1$，其中 $c = \cos 20^\circ$；三次
$8X^3 - 6X - 1$ 无有理根（候选 $\pm1, \pm\frac
12, \pm\frac14, \pm\frac18$ 均不是根），故不可约，次数为
$3$。任一 $60^\circ$ 角都可构造，故三等分者
将构造出 $c$。(3) 若 $\sqrt\pi$ 可构造，则它
为代数数，从而 $\pi$ 也是代数数。关于 $n$ 边形，
$\zeta_n$ 的次数为 $\varphi(n)$
（\cref{thm:b3:galois:cyclotomicirred}）；必要性由
万策尔定理推出；充分性方面，\omterm{def:b3:galois:galois}{伽罗瓦群}为阿贝尔群，阶为
$2^m$，有指数为 $2$ 的子群链（有限 $2$-群
均有此性质，见\cref{exo:b3:groups:10}），其不动域构成
一个二次塔，终于 $\Q(\zeta_n)$
（\cref{thm:b3:galois:fundamental}）；由
\cref{thm:b3:galois:wantzel} 即得结论。
\end{proof}

\section{根式可解性}

\begin{definition}\label{def:b3:galois:radical}
设扩张 $L/K$（本节中特征恒为 $0$），
称其为\emph{根式的}，若存在域塔 $K = F_0 \subseteq \dots
\subseteq F_r = L$，使得 $F_{i+1} = F_i(\alpha_i)$，
$\alpha_i^{n_i} \in F_i$：每一步添加一个 $n_i$ 次方根。
称多项式 $P \in K[X]$ 为\emph{根式可解
}\index{根式可解}，若其\omterm{thm:b3:galois:splitting}{分裂域}
含于 $K$ 的某个根式扩张中。
\end{definition}

\begin{lemma}\label{lem:b3:galois:cyclicsteps}
设域 $K$ 含本原 $n$ 次单位根 $\zeta$，
即 $\zeta$ 的阶为 $n$（于 $K^\times$ 中），且 $a \in K^\times$。则
$K(\sqrt[n]a)/K$ 是\omterm{def:b3:galois:galois}{伽罗瓦扩张}，其\omterm{def:b3:galois:galois}{伽罗瓦群}为\emph{循环群}。反之
——下文不需要——，每个次数为 $n$ 的循环扩张
都具有这种形式。此外，$K(\zeta_n)/K$ 是\omterm{def:b3:galois:galois}{伽罗瓦扩张}，
其\omterm{def:b3:galois:galois}{伽罗瓦群}为\emph{交换群}，对任意 $K$（其特征为 $0$）成立。
\end{lemma}

\begin{proof}
$X^n - a$ 可分（$\gcd$ 与 $nX^{n-1}$ 互素，因 $a \neq 0$），且
在 $K(\alpha)$ 中完全分裂，$\alpha^n = a$，其根为
$\zeta^k\alpha \in K(\alpha)$。于是 $K(\alpha)/K$ 为\omterm{def:b3:galois:galois}{伽罗瓦扩张}；
映射 $\sigma \mapsto \sigma(\alpha)/\alpha \in \mu_n = \langle
\zeta\rangle$ 为单同态（$\sigma\tau(\alpha) =
\sigma(\tau(\alpha)/\alpha \cdot \alpha) =
\tau(\alpha)/\alpha\cdot\sigma(\alpha)$，其中商属于
$K$），映入一个循环群，故 $\operatorname{Gal}$ 为循环群。
逆命题即库默尔理论，本书不需要（见
下文注记）。对于 $K(\zeta_n)$，\omterm{def:b3:galois:separable}{可分多项式} $X^n -
1$ 在其上完全分裂，且 $\sigma \mapsto a(\sigma)$，其中 $\sigma(\zeta_n) =
\zeta_n^{a(\sigma)}$，此映射将该群嵌入交换群
$(\Z/n\Z)^\times$ 中，如\cref{thm:b3:galois:cyclotomicirred}所述
（单射性只需 $\zeta_n$ 生成所涉及的
单位根即可）。
\end{proof}

\begin{theorem}[伽罗瓦]\label{thm:b3:galois:solvable}
设 $K$ 的特征为 $0$，$P \in K[X]$ 的
\omterm{thm:b3:galois:splitting}{分裂域}为 $L$。若 $P$ \omterm{def:b3:galois:radical}{根式可解}，则
$\operatorname{Gal}(L/K)$ 为\omterm{def:b3:groups:derived}{可解群}。（逆命题
也成立，本书不需要它。）\index{伽罗瓦可解性
定理}
\end{theorem}

\begin{proof}
第一步：\emph{将根式域塔扩充为伽罗瓦域塔。}设 $L
\subseteq M$，$M/K$ 为根式扩张，根次数分别为 $n_1, \dots,
n_r$ 与 $n = n_1\cdots n_r$。首先添加 $\zeta_n$：
域塔 $K \subseteq K(\zeta_n) \subseteq M(\zeta_n)$ 仍是\omterm{def:b3:galois:radical}{根式的}
（$\zeta_n$ 是单位根：一步根式添加，$\zeta_n^n = 1$），
且第一步之后的各步都在包含所需
单位根的域上进行。接着，将 $M(\zeta_n)$ 替换为
$N$，即所有 $\sigma(M(\zeta_n))$ 的复合域，其中 $\sigma$ 取遍
（有限多个）$K$-嵌入，它们将 $M(\zeta_n)$ 映入某固定
\omterm{def:b3:galois:closure}{代数闭包}中：$N$ 为一组生成元的\omterm{def:b3:galois:algebraic}{极小多项式}之积的
\omterm{thm:b3:galois:splitting}{分裂域}（特征为 $0$：
有限且可分），故 $N/K$ 为\omterm{def:b3:galois:galois}{伽罗瓦扩张}；又 $N$
相对 $K$ 为根式扩张：每个 $\sigma(M(\zeta_n))$ 相对 $K$ 亦为\omterm{def:b3:galois:radical}{根式的}
（将 $\sigma$ 作用于根式域塔即可），且根式扩张的
复合仍为\omterm{def:b3:galois:radical}{根式的}（把域塔拼接起来：若 $F'/K$ 为
\omterm{def:b3:galois:radical}{根式的}，其域塔添加 $\beta_j$，则形如 $F''(\beta_j)$ 的
各步在任何更大底域上仍为\omterm{def:b3:galois:radical}{根式的}）。

第二步：\emph{由伽罗瓦域塔读出可解性。}因此可设
$L \subseteq N$，$N/K$ 为伽罗瓦且\omterm{def:b3:galois:radical}{根式的}，域塔为 $K
\subseteq K(\zeta_n) = E_0 \subseteq E_1 \subseteq \dots
\subseteq E_s = N$，每步为 $E_{i+1} = E_i(\sqrt[n_i]{a_i})$，且
$\zeta_{n_i} \in E_0 \subseteq E_i$。记 $G =
\operatorname{Gal}(N/K)$ 与 $G_i = \operatorname{Gal}(N/E_i)$：
一条递降链 $G \supseteq G_0 \supseteq G_1 \supseteq
\dots \supseteq G_s = \{e\}$。每个 $E_{i+1}/E_i$ 为\omterm{def:b3:galois:galois}{伽罗瓦扩张}，
其\omterm{def:b3:galois:galois}{伽罗瓦群}为循环群（\cref{lem:b3:galois:cyclicsteps}），故由
伽罗瓦基本定理应用于\omterm{def:b3:galois:galois}{伽罗瓦扩张} $N/E_i$
（\cref{thm:b3:galois:fundamental}(3)，所在群为 $G_i$）：
$G_{i+1} \trianglelefteq G_i$，且 $G_i/G_{i+1} \cong
\operatorname{Gal}(E_{i+1}/E_i)$ 为循环群。类似地，$E_0/K$ 为
\omterm{def:b3:galois:galois}{伽罗瓦扩张}，其\omterm{def:b3:galois:galois}{伽罗瓦群} $G/G_0$ 为交换群
（\cref{lem:b3:galois:cyclicsteps}）。该链表明 $G$
可解（\cref{prop:b3:groups:derived}）。最后，
$\operatorname{Gal}(L/K)$ 是 $G$ 的\emph{\omterm{thm:b3:groups:quotient}{商群}}：$L/K$ 为
\omterm{def:b3:galois:galois}{伽罗瓦扩张}（$P$ 在特征为 $0$ 时可分），且限制映射 $G
\to \operatorname{Gal}(L/K)$ 是满射
（\cref{thm:b3:galois:fundamental}(3)，取 $H =
\operatorname{Gal}(N/L)$）；\omterm{def:b3:groups:derived}{可解群}的\omterm{thm:b3:groups:quotient}{商群}
仍可解。
\end{proof}

\begin{corollary}[五次方程的不可解性]
\label{cor:b3:galois:quintic}
存在次数为 $5$ 的 $\Q$ 上的多项式不是
\omterm{def:b3:galois:radical}{根式可解}的：例如 $X^5 - 4X + 2$，其\omterm{def:b3:galois:galois}{伽罗瓦群}为
$S_5$（\cref{exo:b3:galois:11}），这是一个不\omterm{def:b3:groups:derived}{可解群}
（\cref{cor:b3:groups:snnotsolvable}）。
次数 $\geq 5$ 时，不存在根式形式的一般公式。\index{五次方程，
不可解性}
\end{corollary}

\begin{remark}\label{rem:b3:galois:converse}
\cref{thm:b3:galois:solvable}的逆命题——可解的\omterm{def:b3:galois:galois}{伽罗瓦群}
蕴含\omterm{def:b3:galois:radical}{根式可解}——其证明通过
下降\omterm{def:b3:groups:derived}{导出列}，并证明每个循环扩张
（含足够多单位根）都是\omterm{def:b3:galois:radical}{根式的}，借助\emph{拉格朗日
预解式}；这说明为何次数为 $2, 3, 4$ 时存在求根公式：
$S_2, S_3, S_4$ 均为\omterm{def:b3:groups:derived}{可解群}。
（\cref{ex:b3:groups:solvableexamples}）。该结论在
本阶段不加证明地承认；完整的讨论属于硕士课程，但
\cref{exo:b3:galois:8}对三次情形将其具体化。
\end{remark}

\section{练习}

\begin{exercise}[$\star$]\label{exo:b3:galois:1}
证明 $[\Q(\sqrt2, \sqrt3):\Q] = 4$，且 $\Q(\sqrt2 + \sqrt3) =
\Q(\sqrt2, \sqrt3)$，并求
$\sqrt2 + \sqrt3$ 在 $\Q$ 上的\omterm{def:b3:galois:algebraic}{极小多项式}。
\end{exercise}

\begin{exercise}[$\star$]\label{exo:b3:galois:2}
设 $\alpha = \sqrt[3]2$。证明 $\Q(\alpha)/\Q$ 不是
正规的（给出一个嵌入 $\Q(\alpha) \to \C$，使其像不为
$\Q(\alpha)$），求 $L$（$X^3 -
2$ 的\omterm{thm:b3:galois:splitting}{分裂域}），及 $[L:\Q]$，并验证 $\operatorname{Aut}_\Q(\Q(\alpha)) =
\{\mathrm{id}\}$：对于非\omterm{def:b3:galois:galois}{伽罗瓦扩张}，自同构
群可以远小于\omterm{def:b3:galois:extension}{扩张次数}。
\end{exercise}

\begin{exercise}[$\star$]\label{exo:b3:galois:3}
构造 $\mathbb F_9$，使其为 $\mathbb F_3[X]/(X^2+1)$，并求
$\mathbb F_9^\times$ 的一个生成元。列出
次数分别为 $1, 2, 3$ 的（$\mathbb F_2$ 上）首一不可约多项式，并验证
$X^8 - X = X(X+1)(X^3+X+1)(X^3+X^2+1)$ 在 $\mathbb F_2$ 上成立。
\end{exercise}

\begin{exercise}[$\star\star$]\label{exo:b3:galois:4}
(a) 求模 $7$ 与模 $11$ 的所有原根（即
$\mathbb F_7^\times$、$\mathbb F_{11}^\times$ 的生成元）。
(b) 设 $p$ 为奇素数。证明：对 $x \in \mathbb F_p^\times$，x 是平方元
当且仅当 $x^{(p-1)/2} = 1$（欧拉判别法），并由此得到
\cref{pb:b3:rings:1} 中关于 $-1$ 的判别法。
\end{exercise}

\begin{exercise}[$\star\star$]\label{exo:b3:galois:5}
证明 $\mathbb F_{p^m} \cap \mathbb F_{p^n} = \mathbb
F_{p^{\gcd(m,n)}}$ 及 $\mathbb F_{p^m}\mathbb F_{p^n} = \mathbb
F_{p^{\operatorname{lcm}(m,n)}}$，二者均在固定的\omterm{def:b3:galois:closure}{代数
闭包} $\bar{\mathbb F}_p$ 中，并描述
$\operatorname{Gal}(\mathbb F_{p^n}/\mathbb F_{p^m})$，其中 $m
\mid n$。
\end{exercise}

\begin{exercise}[$\star\star$]\label{exo:b3:galois:6}
设 $I_d(q)$ 为
次数为 $d$ 的 $\mathbb F_q$ 上首一不可约多项式的个数。证明
\[
X^{q^n} - X \;=\; \prod_{d \mid n}\ \prod_{P \text{ 首一
不可约，} \deg P = d} P ,
\qquad\text{因此}\qquad
q^n = \sum_{d \mid n} d\, I_d(q).
\]
由此显式求出 $I_1, I_2, I_3, I_4$，并证明 $I_d(q) \geq 1$
对每个 $d$ 成立（从而扩张 $\mathbb F_{q^d}/\mathbb F_q$ 存在，
可取为\omterm{def:b3:rings:ideal}{商环} $\mathbb F_q[X]/(P)$ 的形式，对所有 $d$ 成立）。
\end{exercise}

\begin{exercise}[$\star\star$]\label{exo:b3:galois:7}
求 $\operatorname{Gal}(\Q(\sqrt2,\sqrt3)/\Q)$ 以及
中间域构成的完整格。同样的问题对
$(X^2-2)(X^2-3)(X^2-6)$ 的\omterm{thm:b3:galois:splitting}{分裂域}进行讨论，你
注意到了什么？
\end{exercise}

\begin{exercise}[$\star\star$]\label{exo:b3:galois:8}
（三次方程，借助其群求解）设 $P = X^3 + pX + q \in
\Q[X]$ 为不可约多项式，根为 $x_1, x_2, x_3$，\omterm{thm:b3:galois:splitting}{分裂域}为
$L$。令 $\delta = (x_1 - x_2)(x_1 - x_3)(x_2 - x_3)$，且
$\Delta = \delta^2 = -4p^3 - 27q^2$（此经典恒等式可承认，
或通过展开对称函数予以验证）。
(a) 证明 $\operatorname{Gal}(L/\Q) \cong A_3$ 或 $S_3$，
取决于 $\Delta$ 是否为 $\Q$ 中的平方数。
(b) 取 $j = \zeta_3$，定义拉格朗日预解式 $u = x_1
+ jx_2 + j^2x_3$ 与 $v = x_1 + j^2x_2 + jx_3$。证明 $u^3 + v^3
= -27q$ 且 $uv = -3p$，并解出 $u^3, v^3$，由此导出卡尔达诺
公式。$S_3$ 的可解性在何处被用到？
\end{exercise}

\begin{exercise}[$\star\star$]\label{exo:b3:galois:9}
在 $\Q(\zeta_5)$ 中证明唯一的二次子域为
$\Q(\sqrt5)$，利用高斯和 $\eta_0 = \zeta_5 + \zeta_5^4$，
$\eta_1 = \zeta_5^2 + \zeta_5^3$：计算 $\eta_0 + \eta_1$ 与
$\eta_0\eta_1$，并由此推出 $\cos\frac{2\pi}5 = \frac{\sqrt5 -
1}4$。由此得正五边形可构造。
\end{exercise}

\begin{exercise}[$\star\star\star$]\label{exo:b3:galois:10}
设 $K = \mathbb F_p(S, T)$（两个不定元的有理
函数域），$L = K(S^{1/p}, T^{1/p})$。
(a) 证明 $[L:K] = p^2$，且 $\alpha^p \in K$ 对
每个 $\alpha \in L$ 成立。
(b) 由此推出 $L/K$ \emph{不是}单扩张：不存在本原元
存在——不可分性使
\cref{thm:b3:galois:primitive} 失效。
\end{exercise}

\begin{exercise}[$\star\star\star$]\label{exo:b3:galois:11}
设 $P = X^5 - 4X + 2$，$G$ 为其在 $\Q$ 上的\omterm{def:b3:galois:galois}{伽罗瓦群}，
作用在 $5$ 个根上。
(a) 证明 $P$ 不可约，由此推出 $5 \mid \abs G$，从而
$G$ 含有一个 $5$-轮换（由柯西定理，
\cref{thm:b3:groups:cauchy}）。
(b) 通过研究 $x \mapsto x^5 - 4x +
2$ 的变化情况，证明 $P$ 恰有 $3$ 个实根；由此推出
复共轭限制在 $G$ 中是一个对换。
(c) 证明：$S_5$ 的一个含有对换与
$5$-轮换的子群必为 $S_5$ \emph{（将对换经
轮换的幂共轭）}。由此得 $G \cong S_5$，并结合
\cref{thm:b3:galois:solvable} 知
$P$ 不\omterm{def:b3:galois:radical}{根式可解}。
\end{exercise}

\begin{exercise}[$\star\star\star$]\label{exo:b3:galois:12}
（二面体四次情形）设 $\alpha = \sqrt[4]2$，$L =
\Q(\alpha, \iu)$，即 $X^4 - 2$ 在 $\Q$ 上的\omterm{thm:b3:galois:splitting}{分裂域}。
(a) 证明 $[L : \Q] = 8$，且 $G =
\operatorname{Gal}(L/\Q)$ 由 $\sigma\colon
\alpha \mapsto \iu\alpha,\ \iu \mapsto \iu$ 与复共轭
$\tau$，满足 $\sigma^4 = \tau^2 = e$ 且
$\tau\sigma\tau = \sigma^{-1}$，从而 $G \cong D_4$。
(b) 列出 $D_4$ 的子群格（十个子群），并
将每个子群与其不动域对应，特别验证
$\Q(\sqrt2)$、$\Q(\iu)$、$\Q(\iu\sqrt2)$ 是三个
二次子域，并定位 $\Q(\alpha)$、
$\Q(\iu\alpha)$、$\Q(\sqrt2, \iu)$。
(c) 哪些中间域是 $\Q$ 上的\omterm{def:b3:galois:galois}{伽罗瓦扩张}？请将
答案与 $D_4$ 的\omterm{def:b3:groups:normal}{正规子群}对照，并
解释为何 $\Q(\alpha)/\Q$ 不是\omterm{def:b3:galois:galois}{伽罗瓦扩张}，而 $\Q(\sqrt2)/\Q$
则是。
\end{exercise}

\section{问题：高斯与正 17 边形}

\begin{problem}[{周末问题——正
17 边形的可构造性}]\label{pb:b3:galois:1}
1796 年 3 月 30 日，年仅十九岁的高斯证明了
正 $17$ 边形是\omterm{def:b3:galois:constructible}{可构造的}——这是自古代以来该问题上的
首个进展。用本章的
工具重构他的计算。设 $\zeta = \eu^{2\iu\pi/17}$，$L =
\Q(\zeta)$，$G = \operatorname{Gal}(L/\Q)$。

\medskip
\noindent\textbf{第一部分——群及其子群链。}
\begin{enumerate}
  \item 说明：$[L:\Q] = 16$，$G \cong (\Z/17\Z)^\times$，
        $16$ 阶循环群。验证 $3$ 是
        $(\Z/17\Z)^\times$ 的生成元 \emph{（计算 $3$ 的幂
        模 $17$：$3, 9, 10, 13, 5, 15, 11, 16, \dots$）}。
  \item 设 $\sigma \in G$ 满足 $\sigma(\zeta) = \zeta^3$，且
        $H_k = \langle \sigma^{2^k}\rangle$，其中 $k = 0, \dots,
        4$。证明 $G = H_0 \supset H_1 \supset H_2 \supset
        H_3 \supset H_4 = \{e\}$，且每个指数 $[H_k :
        H_{k+1}] = 2$，且不动域 $\Q = L_0
        \subset L_1 \subset L_2 \subset L_3 \subset L_4 = L$
        构成一个二次扩张的塔。
  \item 由\emph{先验地}得出（利用
        \cref{thm:b3:galois:wantzel}）：$\zeta$——进而正
        $17$ 边形——是\omterm{def:b3:galois:constructible}{可构造的}。本题的其余部分
        将把这一塔显式地写出。
\end{enumerate}

\medskip
\noindent\textbf{第二部分——长为 8 的周期。}
定义\emph{高斯周期}如下：
\[
\eta_0 = \sum_{k \text{ 偶数}} \zeta^{3^k \bmod 17}
= \zeta^{1} + \zeta^{9} + \zeta^{13} + \zeta^{15} + \zeta^{16} +
\zeta^{8} + \zeta^{4} + \zeta^{2},
\qquad
\eta_1 = \sum_{k \text{ 奇数}} \zeta^{3^k \bmod 17}.
\]
\begin{enumerate}[resume]
  \item 证明 $\eta_0, \eta_1$ 被 $H_1$ 固定，
        并被 $\sigma$ 互换；由此推出 $\eta_0, \eta_1 \in L_1$，且
        它们是 $\Q$ 上某一二次方程的两个根。
  \item 计算 $\eta_0 + \eta_1 = -1$，证明 $\eta_0\eta_1 = -4$
        \emph{（每个乘积 $\zeta^a\zeta^b$ 都等于某个 $\zeta^c$，
        $c \neq 0$；数一数每个 $c$ 出现的次数，或
        论证该乘积是被 $G$ 固定的有理整数，
        它等于全部 $64$ 个乘积之和，并利用
        每个非零剩余类出现的次数相同）}。
  \item 推出 $\eta_0 = \frac{-1 + \sqrt{17}}2$，$\eta_1 =
        \frac{-1-\sqrt{17}}2$（数值上辨别哪个是哪个：
        $\eta_0 \approx 1.56$），并且 $L_1 =
        \Q(\sqrt{17})$.
\end{enumerate}

\medskip
\noindent\textbf{第三部分——长为 4 与 2 的周期。}
定义如下：
\[
\beta_0 = \zeta + \zeta^{13} + \zeta^{16} + \zeta^{4},\quad
\beta_1 = \zeta^3 + \zeta^5 + \zeta^{14} + \zeta^{12},\quad
\beta_2 = \zeta^9 + \zeta^{15} + \zeta^{8} + \zeta^{2},\quad
\beta_3 = \zeta^{10} + \zeta^{11} + \zeta^{7} + \zeta^{6}.
\]
\begin{enumerate}[resume]
  \item 证明 $\beta_0 + \beta_2 = \eta_0$，$\beta_1 + \beta_3 =
        \eta_1$，且 $\beta_0, \beta_2$ 被 $H_2$ 固定，
        并被 $\sigma^2$ 互换。
  \item 计算 $\beta_0\beta_2 = -1$ 与 $\beta_1\beta_3 = -1$
        \emph{（展开：所得的十六个指数恰好覆盖
        $1, \dots, 16$ 各一次）}。
  \item 推出 $\beta_0 = \frac{\eta_0 + \sqrt{\eta_0^2 + 4}}2$
        （数值上检验符号：$\beta_0 \approx 2.05$），以及
        $\beta_1$ 的类似公式；从而 $L_2 =
        \Q(\beta_0)$，是 $L_1$ 上的二次扩张。
  \item 设 $\gamma_0 = \zeta + \zeta^{16} =
        2\cos\frac{2\pi}{17}$，$\gamma_1 = \zeta^{13} +
        \zeta^4$。证明 $\gamma_0 + \gamma_1 = \beta_0$，且
        $\gamma_0\gamma_1 = \beta_1$，从而 $\gamma_0 =
        \frac{\beta_0 + \sqrt{\beta_0^2 - 4\beta_1}}2$.
  \item 整理出用逐层平方根表示
        $\cos\frac{2\pi}{17}$ 的公式链，并给出
        小数验算（$\cos\frac{2\pi}{17} \approx 0.93247$）。
\end{enumerate}

\medskip
\noindent\textbf{第四部分——尾声。}
\begin{enumerate}[resume]
  \item 论证中究竟在何处用到了 $17$ 是
        \emph{费马素数}（$17 = 2^{2^2} + 1$）？证明：
        对素数 $p$，正 $p$ 边形可构造当且仅当 $p
        = 2^{2^t} + 1$ 对某个 $t$ 成立 \emph{（若 $p - 1 = 2^m$，
        证明 $m$ 本身必为 $2$ 的幂）}。
  \item 推出全部\omterm{def:b3:galois:constructible}{可构造的}正
        $n$ 边形（$n \leq 20$）的完整列表，利用高斯--万策尔
        判别法（见\cref{cor:b3:galois:impossible}）。
\end{enumerate}

\medskip
\noindent\textbf{第五部分——高斯和与
二次互反律。} 第二部分的周期中隐藏着一件宝藏。对
奇素数 $p$，\emph{勒让德符号} $\bigl(\frac
ap\bigr)$ 的值为 $+1$，若 $a$ 是模 $p$ 的非零平方；为 $-1$，若
它不是；为 $0$，若 $p \mid a$。由\cref{exo:b3:galois:4}(b)
（欧拉判别法）给出 $\bigl(\frac ap\bigr) \equiv
a^{(p-1)/2} \pmod p$，由此得乘性。记 $\zeta =
\eu^{2\iu\pi/p}$，$p^* = (-1)^{(p-1)/2}p$，并定义
\emph{高斯和}如下：
\[
g \;=\; \sum_{a=1}^{p-1}\Bigl(\frac ap\Bigr)\zeta^a .
\]
\begin{enumerate}[resume]
  \item 证明 $\sum_{a=1}^{p-1}\bigl(\frac ap\bigr) = 0$
        （平方数与非平方数一样多），并证明
        另一形式 $g = \sum_{a=0}^{p-1}\zeta^{a^2}$
        \emph{（每个非零平方恰被取到两次，并且
        $\sum_{a}\zeta^a = 0$）}。对 $p = 17$，将 $g$ 与
        第二部分的周期相联系——证明 $g = \eta_0 - \eta_1$
        \emph{（模 $17$ 的平方数恰为生成元 $3$ 的
        偶数次幂）}。
  \item 证明 $g^2 = p^*$：展开
        \[
        g^2 = \sum_{a,b\neq0}\Bigl(\frac{ab}p\Bigr)
        \zeta^{a+b}
        = \sum_{c}\ \sum_{a \neq 0}\Bigl(\frac{a(c -
        a)}p\Bigr)\zeta^{c}
        \]
        （令 $b = c - a$），代入 $c - a = at$ 以计算
        内和为 $\bigl(\frac{-1}p\bigr)(p - 1)$（当
        $c = 0$ 时），其余情形为 $-\bigl(\frac{-1}p\bigr)$，并
        由第 14 问得出结论。数值验证：对 $p =
        17$，$(\eta_0 - \eta_1)^2 = 17$（见第二部分）。
  \item 推出 $\sqrt{p^*} \in \Q(\zeta_p)$，并得出
        \emph{唯一}的二次子域
        在 $\Q(\zeta_p)$ 中为 $\Q(\sqrt{p^*})$——唯一性
        源于 $\operatorname{Gal}(\Q(\zeta_p)/\Q)$ 为循环群，
        （见\cref{thm:b3:galois:cyclotomicirred}），而循环群
        恰有一个指数为 $2$ 的子群。（每个二次域都可嵌入
        某个分圆域——
        这是克罗内克--韦伯定理的第一种情形，
        其一般形式将在后文很远处才出现。）
  \item 现设 $q \neq p$ 为另一奇素数。在
        环 $\Z[\zeta]$ 中模 $q$ 进行计算，证明
        \[
        g^q \equiv \Bigl(\frac qp\Bigr)\,g \pmod{q\Z[\zeta]}
        \]
        \emph{（新生梦：$(x + y)^q \equiv x^q + y^q$
        在任意交换环中模 $q$ 成立；于是 $g^q \equiv
        \sum_a\bigl(\frac ap\bigr)^q\zeta^{aq}$，作代换 $b =
        aq$，并提出 $\bigl(\frac{q^{-1}}p\bigr) =
        \bigl(\frac qp\bigr)$)}.
  \item 另一方面，$g^q = g\,(g^2)^{(q-1)/2} =
        g\,(p^*)^{(q-1)/2}$。利用模 $q$ 的欧拉判别法，
        推出 $g^q \equiv \bigl(\frac{p^*}q\bigr)g
        \pmod{q\Z[\zeta]}$。然后将关于 $g^q$ 的两个表达式
        各乘以 $g$，并利用 $g^2 = p^*$，
        它模 $q$ 可逆——于是得出
        \[
        \Bigl(\frac qp\Bigr) = \Bigl(\frac{p^*}q\Bigr) .
        \]
        \emph{（为何整数之间的
        同余 $\pm p^*$（模 $q\Z[\zeta]$）能蕴含二者相等？
        与 $\Z$ 求交即可。）}。
  \item 展开 $\bigl(\frac{p^*}q\bigr) =
        \bigl(\frac{-1}q\bigr)^{(p-1)/2}\bigl(\frac pq\bigr)$
        及 $\bigl(\frac{-1}q\bigr) = (-1)^{(q-1)/2}$，以
        得到\emph{二次互反律}：
        \[
        \Bigl(\frac pq\Bigr)\Bigl(\frac qp\Bigr)
        = (-1)^{\frac{p-1}2\cdot\frac{q-1}2} .
        \]
        以 $(p, q) = (17, 3)$ 为例验证，列出
        模 $17$ 与模 $3$ 的平方数，并借此在三行内
        判断 $x^2 \equiv 219 \pmod{383}$ 是否可解
        （$383$ 为质数，$219 = 3\cdot73$）。
\end{enumerate}

\medskip
\noindent\textbf{第六部分——计算不可约多项式的个数：
$\mathbb F_q[X]$ 的质数定理。} 固定一个质数幂
       $q$，并设 $N_q(n)$ 为次数为 $n$ 的\emph{首一不可约}
       多项式在 $\mathbb F_q$ 上的个数。回顾
       \cref{exo:b3:galois:6} 中 $X^{q^n} - X$ 的因式分解，
       及恒等式 $q^n = \sum_{d\mid n}d\,N_q(d)$，
       现将其求逆、重新诠释并加以利用。
\begin{enumerate}[resume]
  \item （字）称字 $w \in \mathbb F_q^n$
        为\emph{本原}的，若它不是某个严格更短的字的幂 $u^{n/d} =
        u\cdots u$（$u$ 为严格更短的字），并设
        $A(d)$ 为长度为 $d$ 的本原字的个数。
        证明：长度为 $n$ 的每个字都唯一地是某个本原字的幂，
        该本原字的长度为 $d \mid n$，因而
        $q^n = \sum_{d \mid n}A(d)$；与
        \cref{exo:b3:galois:6} 比较，得出 $A(d) = d\,N_q(d)$
        对所有 $d$ 成立，并以一个显式双射解释这一巧合：
        设 $\alpha \in
        \mathbb F_{q^d}$ 为次数为 $d$ 的元素，其弗罗贝尼乌斯\omterm{def:b3:groups:action}{轨道}
        $(\alpha, \alpha^q, \dots, \alpha^{q^{d-1}})$ 恰含
        $d$ 个互异元素。次数为 $d$ 的元素与
        次数为 $d$ 的不可约多项式之间
        成 $d$ 对一的对应。
  \item 证明\emph{莫比乌斯反演公式}：若 $f(n) =
        \sum_{d\mid n}g(d)$ 对所有 $n$ 成立，则 $g(n) =
        \sum_{d\mid n}\mu(d)\,f(n/d)$，其中 $\mu$ 为
        莫比乌斯函数（$\mu(m) = (-1)^{\#\text{素因子个数}}$，
        若 $m$ 无平方因子，否则为 $0$）
        \emph{（关键引理：$\sum_{d \mid m}\mu(d) = 0$ 对 $m >
        1$ 成立 —— 将因数按是否含某一固定素因子配对）}。
        由此推得
        \[
        N_q(n) = \frac1n\sum_{d \mid n}\mu(d)\,q^{n/d} .
        \]
  \item 证明：$N_q(n) \geq \frac1n\bigl(q^n - 2q^{n/2}\bigr) > 0$
        对每个 $n \geq 1$ 成立，这给出 $\mathbb F_{q^n}$
        对所有 $n$ 存在的一个新证明。解释首项的含义：
        随机取一个次数为 $n$ 的首一多项式，它不可约的
        概率 $\sim \frac1n$，这正是质数定理的完美类比，
        其中 $\log x$ 换成了 $n$；
        对 $q = 2$、$n \leq 4$ 作数值验证
        （个数见\cref{exo:b3:galois:6}）。
  \item 证明第 21 问的乘法版本：
        \[
        \prod_{\substack{\pi \text{ 首一不可约}\\ \deg\pi =
        n}}\pi \;=\; \prod_{d \mid n}
        \bigl(X^{q^d} - X\bigr)^{\mu(n/d)}
        \]
        \emph{（非零有理函数构成的阿贝尔群上的
        莫比乌斯反演）}；对 $q
        = 2$、$n = 2$ 手算验证：$(X^4 - X)/(X^2 - X) = X^2 + X + 1$。
\end{enumerate}

\medskip
\noindent\textbf{第七部分——两个尾声。}
\begin{enumerate}[resume]
  \item （第二补充律）第五部分的方法也可计算
        $\bigl(\frac2q\bigr)$。设 $\omega = \eu^{2\iu\pi/8}$
        并令 $g = \omega + \omega^{-1}$，证明 $g^2 = 2$
        \emph{（$\omega^2 = \iu$）}；然后对奇质数 $q$，
        证明在 $\Z[\omega]$ 中模 $q$ 下式成立：
        \[
        g^q \equiv \omega^q + \omega^{-q} \pmod{q\Z[\omega]},
        \]
        并证明右端等于 $g$（当 $q \equiv \pm1
        \pmod 8$ 时），等于 $-g$（当 $q \equiv \pm3 \pmod 8$ 时）。将此与
        $g^q = g\,(g^2)^{(q-1)/2} \equiv
        \bigl(\frac2q\bigr)g$ 比较（参见第 18 问），得出
        \[
        \Bigl(\frac2q\Bigr) = (-1)^{(q^2-1)/8},
        \]
        并核对 $(q^2 - 1)/8$ 为偶数当且仅当 $q
        \equiv \pm1 \pmod 8$。再验证：$2$ 是模 $7$ 的平方，
        模 $17$ 亦然（$3^2$ 与 $6^2$），而模 $3$ 与
        模 $5$ 则否。
  \item （$\mathbb F_q[X]$ 的泽塔函数）证明形式幂级数
        关于 $t$ 的恒等式：
        \[
        \prod_{n \geq 1}\bigl(1 - t^n\bigr)^{-N_q(n)}
        = \frac1{1 - qt}
        \]
        \emph{（唯一分解为首一不可约多项式：
        将每个因子展开为几何级数，并计数
        次数为 $n$ 的首一多项式）}。通过取对数，恢复
        $q^m = \sum_{d \mid m}d\,N_q(d)$ 这一恒等式
        （见\cref{exo:b3:galois:6}）。
        手算 $t^2$ 的系数，取 $q = 2$，并利用
        第 21 问的公式计算 $N_2(6) = 9$，
        并验证 $2^6 = 1\cdot2 + 2\cdot1 + 3\cdot2 +
        6\cdot9$。
\end{enumerate}
\end{problem}
